\documentclass[critical]{pol-ren}
\usepackage{gregoriotex}
%% Text of stanzas 1–2, shared by both editions, and the local macros.

%% Diplomatic text of K (back flyleaf of BJ MS 1619).
\newcommand\diplOne{Bogv rodzicza dzewicza bogem slawena maria\\
  Vtwego syna gospodzina matko swolena maria\\
  Sziszczi nam spwczi nam kyrieleyson}
\newcommand\diplTwo{Twego dzela krzcziczela boszicze\\
  Vslisz glosi naplen misli czlowecze\\
  Slisz modlitwǫ yǫsz nosimi\\
  A dacz raczi gegosz prosimi\\
  a naswecze zbozni pobith\\
  posziwocze raski pzbith ky'eleyson}

%% Critical text, normalised after Woronczak (1962).
\newcommand\stanzaOne{Bogurodzica dziewica, Bogiem sławiena Maryja,\\
  u twego syna Gospodzina matko zwolena, Maryja!\\
  Zyszczy nam, spuści nam. Kirielejson.}
\newcommand\stanzaTwo{Twego dziela Krzciciela, Bożyce,\\
  usłysz głosy, napełń myśli człowiecze.\\
  Słysz modlitwę, jąż nosimy,\\
  a dać raczy, jegoż prosimy:\\
  a na świecie zbożny pobyt,\\
  po żywocie rajski przebyt. Kirielejson.}

%% Translation, line by line with the critical text.
\newcommand\transOne{Mother of God, Virgin glorified by God, Mary!\\
  Chosen Mother, Mary, from thy Son the Lord\\
  win for us, send down to us. Lord, have mercy.}
\newcommand\transTwo{For thy Baptist's sake, Son of God,\\
  hear our voices, fulfil the thoughts of men.\\
  Hear the prayer we bring,\\
  and deign to grant what we ask:\\
  on earth a godly life,\\
  after life a dwelling in paradise. Lord, have mercy.}

%% Performance edition: the text as sung, line by line with its translation.
\newcommand\sungOne{Bogurodzica dziewica, Bogiem sławiena Maryja,\\
  u twego syna Gospodzina matko zwolena, Maryja!\\
  Zyszczy nam, spuści nam.\\
  Kyrie eleison.}
\newcommand\sungTwo{Twego dziela Krzciciela, Bożyce,\\
  usłysz głosy, napełń myśli człowiecze.\\
  Słysz modlitwę, jąż nosimy,\\
  a dać raczy, jegoż prosimy:\\
  a na świecie zbożny pobyt,\\
  po żywocie rajski przebyt.\\
  Kyrie eleison.}
\newcommand\sungTransOne{Mother of God, Virgin glorified by God, Mary!\\
  Chosen Mother, Mary, from thy Son the Lord\\
  win for us, send down to us.\\
  Lord, have mercy.}
\newcommand\sungTransTwo{For thy Baptist's sake, Son of God,\\
  hear our voices, fulfil the thoughts of men.\\
  Hear the prayer we bring,\\
  and deign to grant what we ask:\\
  on earth a godly life,\\
  after life a dwelling in paradise.\\
  Lord, have mercy.}




\seriesno{1}
\shortcomposer{Anonymous}
\title{Bogurodzica}
\forces{the two oldest stanzas with their melody}
\sourceline{from Kraków, Biblioteka Jagiellońska, MS 1619}
\edition{First edition}{2026}
\slug{bogurodzica}

\begin{document}
\maketitlepage

\colophon

%% ================================================================ PREFACE
\section*{Introduction}

\initial{B}{ogurodzica} is the oldest surviving Polish religious song. Its
first two stanzas, addressed to the Virgin and to Christ, form the original
song; later copies add up to twenty more. The only medieval copy with music
for both stanzas is a flyleaf of BJ MS 1619, written shortly after 1400. This
edition gives that melody in square notation, with the text of the same copy,
an apparatus of musical and textual variants, and notes on the manuscript.

\subsection{Form} Each stanza has three parts: two parallel lines, then a
contrasting third with the refrain (Łoś). The melody varies four short
phrases (Surzyński's \textit{a–d}), all traceable to the first line, which
turns almost symmetrically on c′–d′–c′ (Flotzinger; Perz calls the
procedure centonate). The cadence of \pl{Maryja} returns at \pl{zwolena},
the second \pl{Maryja} and \pl{Bożyce}. \pl{Dziewica} returns at \pl{zyszczy
nam}, \pl{usłysz głosy} and \pl{a dać raczy}. \pl{Twego dziela Krzciciela}
returns at \pl{napełń myśli człowiecze}, and the end of \pl{a na świecie
zbożny pobyt} at \pl{po żywocie rajski przebyt}.

\subsection{Mode} The final is D, with inner cadences on a and F. Chybiński
weighed mode~6 because of the F cadences, but settled on mode~1. B appears
once, at the melodic peak on \pl{spuści}.

\subsection{Date and origin} Attribution to St Adalbert (d.\,997) begins with
15th- and 16th-century copies, Łaski among them, and no modern scholar
accepts it. Łoś placed the text in the late 13th century on linguistic and
metrical grounds; Chybiński and Feicht put the melody at the turn of the 13th
and 14th centuries. Woronczak (1952) matched the opening phrase to the Holy
Saturday litany \la{Kyrie} in the Graz gradual MS 807 and read the song as
a Kyrie trope. Flotzinger (2005) showed that the match covers only the first
phrase, placed the Graz manuscript at Passau c.\,1170, and rejected the trope
reading. He dates the song to the late 14th century, in south-eastern Poland
with Czech mediation. Feicht (1962) assembled many parallels, among them
Jehan de Braine's \textit{Par desous l'ombre d'un bois}, and judged the melody
the coherent work of one trained musician. Pikulik (2005) argues that the
octave leap after the first \pl{Maryja} rules out a date before the 12th or
13th century and calls the melody a 13th-century cento. Kalina (1880) heard
two melodies by different authors; Feicht's analysis superseded this.

\section*{Sources}

\litgroup{Musical witnesses}
\begin{sigla}
\siglum{K} Kraków, Biblioteka Jagiellońska, MS 1619, back flyleaf.
  \textit{Base text.} Added to a copy of Latin sermons made by Maciej of
  Grochów, vicar of Kcynia; before 1408. Stanzas 1–2 with music, in seven
  systems. Gothic chant notation of German–Messine type, C and F clefs, no
  custodes. A decorated initial opens each musical section: \pl{B(ogu)},
  \pl{U (twego)}, \pl{Z(yszczy)}, \pl{T(wego)}, \pl{U(słysz)}, \pl{S(łysz)},
  \pl{A (dać)}.
\siglum{K70} Kraków, Biblioteka Jagiellońska, manuscript of 1570.
  Leans towards b\fl\ before cadences on~a.
\siglum{W} Warsaw, library of the Literati Chapel at St John's,
  MS 998\,D. Mensural notation, the only witness with written rhythm.
\siglum{Cz} Częstochowa, Jasna Góra, late 15th century. Lost; known
  from a drawing reproduced by Woronczak. Music for the first two lines of
  stanza~1 only.
\siglum{Gn} The Gniezno cathedral melody, printed in \textit{Przyjaciel
  Ludu} 4, no.\,46 (Leszno, 1838), p.~364. It sets the whole text, with
  smaller note values and sharps.
\siglum{Z} Warsaw, Zamoyski Library, MS 1152 (1672). Mensural; the
  \sig{K70} version, less accurately written.
\end{sigla}

\noindent Readings of \sig{K70} and \sig{W} are given after Chybiński (1907,
pp.~30–41), who sets them out in modern notation. \sig{Cz}, \sig{Gn} and
\sig{Z} are described but not collated.

\litgroup{Text witnesses}
\begin{sigla}
\siglum{K408} Kraków, BJ MS 408, f.~87\textsuperscript{v}. Added in the early 15th century to
  a codex dated 1408. Stanzas 1–14, no music.
\siglum{Wa} Warsaw, Biblioteka Narodowa, from Holy Cross, c.\,1456.
  19 stanzas.
\siglum{Cz} As above.
\siglum{K3} Kraków, first half of the 16th century. 23 stanzas.
\siglum{Ł} Jan Łaski, \la{Commune incliti Poloniae Regni privilegium}
  (Kraków: J.\,Haller, 1506), first in the \la{Registrum}, attributed to
  St Adalbert. No music.
\siglum{M} Mateusz z Kościana, 1543.
\siglum{H} Herburt, 1570.
\siglum{S} Skarga, \pl{Żywoty świętych}, 1579.
\end{sigla}

\noindent Pilat (1879) and Kalina (1880) number the Kraków copies the other
way round, calling BJ 408 ‘Kraków I’. The sigla here follow the shelfmarks.

\section*{Editorial method}

\subsection{Melody and text} The melody and text are those of \sig{K}. The
transcription was made from the facsimile and agrees in every pitch with
Chybiński's (1907).

\subsection{Pitch} Pitch is fixed by the clefs of \sig{K}: final D, range
C–d′, mode~1. The single b, on \pl{spuś-ci}, is natural.

\subsection{Rhythm} \sig{K} has no rhythmic signs, so the score has none.

\subsection{Text} The text is normalised after Woronczak (1962), keeping the
\pl{u} that \sig{K} writes and sets to music at the start of line 1.2.

\subsection{Bar lines} Bar lines mark the verse structure.

%% ================================================================ TEXT
\opensection{Text and translation}

\begin{textcols}
& \colhead{critical text} & \colhead{translation}\\[0.6em]
\stanzanumber{1} & \lines{\stanzaOne} & \lines{\transOne}\\[1.1em]
\stanzanumber{2} & \lines{\stanzaTwo} & \lines{\transTwo}\\
\end{textcols}

\section*{Diplomatic text}

\noindent The diplomatic text is that of \sig{K}; \dipl{ǫ} is a nasal vowel.

\begin{textcols}
\stanzanumber{1} & \lines{\diplOne} & \\[1.1em]
\stanzanumber{2} & \lines{\diplTwo} & \\
\end{textcols}

%% ================================================================ SCORE
\scorehead
\gregorioscore[a]{bogurodzica}
\musicnote{Bar lines are editorial. Neume groups follow \sig{K}.}

%% ================================================================ NOTES
\opensection{Critical notes}

\section*{Musical variants}

\noindent Pitches C D E F G a b c′ d′; hyphens join notes on one syllable.
(?!) marks readings Chybiński thought corrupt.

\begin{variants}
\cl{1.1}{Bogurodzica}{\sig{K} D C F E-D-C D \quad \sig{K70} D D-C F E-C-E D \quad \sig{W} D C F E D}
\cl{1.2}{dziewica}{\sig{K} F G-F-G a \quad \sig{K70} F G b\fl\ a \quad \sig{W} F G a}
\cl{1.3}{Bogiem sławiena}{\sig{K} c′ d′ c′ a-G-F F \quad \sig{K70} b d′ c′ a-G F (?!) \quad \sig{W} c′ d′ c′ a F}
\cl{1.4}{Maryja}{\sig{K} F-G-a-G E-C D \quad \sig{K70} F-a-G-E C E D (?!) \quad \sig{W} F a G E D D (?!)}
\cl{1.5}{U twego syna}{\sig{K} d′ c′ a-G-F G-F-G a \quad \sig{K70} a G F G b\fl\ a a \quad \sig{W} short form beginning on a; \sig{K70} and \sig{W} both lack d′ c′}
\cl{1.6}{Gospodzina}{\sig{K} c′ d′ c′ a G F \quad \sig{K70} c′ d′ c′ a-G F \quad \sig{W} c′ d′ c′ a F}
\cl{1.7}{Matko}{\sig{K} c′ a-G-F \quad \sig{K70} c′ a-G (?!) \quad \sig{W} c′ a}
\cl{1.8}{zwolena, Maryja}{as 1.4 in all witnesses}
\cl{1.9}{Zyszczy nam}{\sig{K} F G-F-G a \quad \sig{K70} F G b\fl\ a \quad \sig{W} F G a}
\cl{1.10}{spuści nam}{\sig{K} \sig{W} c′ b d′ \quad \sig{K70} c′ b d′ d′ (?!)}
\cl{1.11}{Kyrie}{\sig{K} c′ a G-F \quad \sig{K70} \sig{W} c′ a F}
\clbreak
\cl{2.1}{Twego dziela Krzciciela}{\sig{K} D-E D-C F G E C D \quad \sig{K70} D D-C F F G G E E C D D (?!) \quad \sig{W} D D-C F G E C-E D}
\cl{2.2}{Bożyce}{as 1.4 in all witnesses}
\cl{2.3}{usłysz głosy}{\sig{K} a G-F G-F-G a \quad \sig{K70} a G-F G b\fl\ a \quad \sig{W} a G-F G a a}
\cl{2.4}{napełń myśli człowiecze}{\sig{K} D C F G E C D \quad \sig{K70} \sig{W} \pl{napełnij}, with an extra note}
\cl{2.5}{Słysz modlitwę, jąż nosimy}{\sig{K} C F G E-D-E C D E D \quad \sig{K70} C G a b\fl\ a a \,|\, D D C C F F E D \quad \sig{W} C F G a \,|\, D D C F G a a}
\cl{2.6}{a dać raczy}{\sig{K} a G-F G-F-G a \quad \sig{K70} a G F G b\fl\ a a (?!)}
\cl{2.7}{jegoż prosimy}{\sig{K} D-E D-C F G a \quad \sig{K70} D D-C F G b\fl\ a}
\cl{2.8}{a na świecie zbożny pobyt}{\sig{K} c′ a c′ G c′ a-G-F G-F-G a \quad \sig{K70} c′ b c′ d′ G G c′ a G G b\fl\ a \quad \sig{W} a c′ d′ G G c′ a F G a a}
\cl{2.9}{po żywocie rajski przebyt}{\sig{K} a c′ d′ G c′ a-G-F G-F-G a \quad \sig{K70} a b\fl\ c′ d′ G \,|\, c′ a F F b\fl\ a \quad \sig{W} a c′ d′ G c′ a F G a a}
\cl{2.10}{Kyrieleison}{\sig{K} a G F F-G-a-G E D D}
\end{variants}

\section*{Textual variants}

\noindent Variants are normalised; spelling differences are ignored.

\begin{variants}
\cl{1.1}{Bogu rodzica}{\sig{K} \sig{K408} \quad \pl{Boga rodzica} \sig{Wa} \sig{Cz} \sig{K3} \sig{Ł} \sig{M} \sig{H} \sig{S}}
\cl{1.1}{sławiena}{\sig{K} \sig{K408} \quad \pl{sławiona} \sig{Wa} \sig{Cz} \sig{K3} \sig{Ł} \sig{M} \sig{H} \sig{S}}
\cl{1.2}{U twego}{\sig{K} \sig{Wa} \sig{Cz} \sig{K3} \sig{Ł} \sig{H} \sig{S} \quad \pl{W twego} \sig{K408} \sig{M}}
\cl{1.2}{Gospodzina}{\sig{K} \sig{K408} \sig{Wa} \sig{Cz} \sig{Ł} \sig{M} \quad \pl{gospodina} \sig{K3} \quad \pl{Gospodynia} \sig{H} \quad \pl{hospodyna} \sig{S}}
\cl{1.2}{zwolena}{\sig{K} \sig{K408} \quad \pl{zwolona} \sig{Wa} \sig{Ł} \sig{M} \sig{H} \sig{S} \quad \pl{zbolenia} \sig{Cz} \quad \pl{zwolienia} \sig{K3}}
\cl{1.3}{Zyszczy \dots\ spuści nam}{\sig{K} \sig{K408} \sig{Wa} \sig{Ł} \quad \pl{zyszczysz \dots\ spuszczysz} \sig{Cz} \quad \pl{zyszczy \dots\ spust winam} \sig{K3} \sig{M} \quad \pl{ziścisz \dots\ spuścisz} \sig{H} \quad \pl{ziści \dots\ spuści} \sig{S}}
\clbreak
\cl{2.1}{Twego dziela}{\sig{K} \sig{K408} \sig{Wa} \sig{Cz} \quad \pl{Twego syna} \sig{K3} \sig{Ł} \sig{M} \sig{H} \sig{S}}
\cl{2.1}{Krzciciela}{\pl{zbawiciela} \sig{M} \quad \pl{[Kr]ziciela zbawiciela} \sig{K3}}
\cl{2.1}{Bożyce}{\sig{K} \sig{K408} \quad \pl{zbożnica} \sig{Wa} \sig{Cz} \quad \pl{zbożny czas} \sig{Ł} \sig{H} \sig{S} \quad \pl{zbożnika} \sig{M}}
\cl{2.2}{napełń}{\sig{K} \quad \pl{napełni} \sig{K408} \sig{Wa} \sig{Cz} \sig{K3} \sig{Ł} \sig{M} \sig{H} \sig{S}}
\cl{2.3}{jąż nosimy}{\sig{K} \sig{K408} \sig{Wa} \quad \pl{jenżeć nosimy} \sig{Cz} \quad \pl{jenże (cię) prosimy} \sig{Ł} \sig{M} \sig{H} \sig{S}}
\cl{2.4}{A dać raczy}{\sig{K} \sig{Wa} \quad \pl{Oddać raczy} \sig{K408} \sig{M} \quad \pl{O dać} \sig{Cz} \quad \pl{O o dać} \sig{Ł} \quad \pl{O o daci} \sig{H} \quad \pl{To dać} \sig{S}}
\cl{2.5}{a na świecie}{\sig{K} \sig{K408} \quad \pl{daj na świecie} \sig{Wa} \sig{Cz} \sig{Ł} \sig{H} \sig{S} \quad \pl{daj nam na świecie} \sig{M} \quad \pl{dai na świecie dobry} \sig{K3}}
\cl{2.5–6}{pobyt \dots\ przebyt}{\pl{pobytk \dots\ przybytk} \sig{Wa}}
\cl{2.6}{rajski}{\dipl{raski} \sig{K}}
\cl{refr.}{Kirielejson}{\dipl{kyrieleyson} \sig{K} (st.\,2 \dipl{ky'eleyson}) \quad \dipl{kyon} \sig{Cz} \quad \dipl{kyeon} \sig{Wa} \quad \pl{Kyrie eleison} \sig{M} \sig{H} \sig{S}}
\end{variants}

\section*{Notes}

\subsection{Staves and clefs} System 1 has four lines, the rest five; the
lowest line also serves as the text baseline. In system~2 the clefs move for
\pl{Matko} and return before \pl{zwolena}. Read with these clefs,
\pl{Matko} is c′ a-G-F.

\subsection{1.1} \dipl{Bogv rodzicza} is two words, with the dative
\pl{Bogu} in a genitive sense (Łoś). All other copies have \pl{Boga
rodzica}, the name used in the 1410 \la{Chronica conflictus} and by Długosz.
\pl{Sławiena} and \pl{zwolena} are Old Polish forms, not Czechisms. The C on
\pl{-gu-} is faint but legible.

\subsection{1.2} \pl{U}: two separate puncta, d′ c′, carry the syllable. The
later musical witnesses drop both notes and so lose the octave leap from
\pl{Maryja}. Over \pl{-dzi-na} of \pl{Gospodzina} the a and G stand close
and the F apart, which gives c′ a-G \,|\, F.

\subsection{2.1} \pl{dziela} means ‘for the sake of’ (later \pl{dla}). Łoś
knew it in no other Old Polish text and used it to date the song to the 13th
century. The printed tradition replaced it with \pl{syna}; modern popular
texts misread it as \pl{dzieła}. \pl{Bożyce} is the vocative of \pl{bożyc},
‘son of God’.

\subsection{2.2} The C on \pl{-pełń} is not visible in the facsimile; only
the D over \dipl{naplen} can be read. It is supplied from the parallel 2.1
and from every other transcription.

\subsection{2.4} The initial of \pl{A dać} has the same design as the \pl{U}
of \pl{Usłysz}. This likely explains Pilat's reading \dipl{O dacz}; Łoś,
Kalina and Woronczak read \dipl{A}.

\subsection{2.4–10} From \pl{A dać raczy} the neumes are larger and more
widely spaced, possibly in another hand (Chybiński). The forms and pitches
are the same as before.

\section*{Literature}

\begin{literature}
\lit Adamiec, M., ed. \textit{Bogurodzica}. Wirtualna Biblioteka Literatury
  Polskiej, Uniwersytet Gdański. \url{https://literat.ug.edu.pl/bgrdzc/}.
\lit Chmiel, A. ‘Uwagi archiwalno-paleograficzne nad pieśnią
  \textit{Bogurodzica} w rękopisie Biblioteki Jagiell. nr 1619’.
  \textit{Rozprawy Wydziału Filologicznego AU} 40 (1904).
\lit Chybiński, A. \textit{‘Bogurodzica’ pod względem historyczno-muzycznym}.
  Kraków, 1907.
\lit Feicht, H. \textit{Studia nad muzyką polskiego średniowiecza}. Opera
  musicologica 1. Kraków: PWM, 1975.
\lit Flotzinger, R. ‘Zur \textit{Bogurodzica}-Frage’. \textit{Context of
  Musicology} 1 (1997): 15–19. Polish trans. \textit{Pamiętnik Literacki}
  96/2 (2005): 7–10.
\lit Gołos, J. Review of Woronczak et al., \textit{Bogurodzica}. \textit{The
  Polish Review} 7/4 (1962): 84–85.
\lit Kalina, A. \textit{Rozbiór krytyczny pieśni ‘Bogarodzica’}. Lwów, 1880.
\lit Łoś, J., ed. \textit{Bogurodzica. Pierwszy polski hymn narodowy}.
  Lublin, 1922.
\lit Perz, M. ‘Polskie “posłowie” do uwag Rudolfa Flotzingera’.
  \textit{Pamiętnik Literacki} 96/2 (2005): 11.
\lit Pikulik, J. ‘Co melodia \textit{Bogurodzicy} mówi nam o czasie
  powstania pieśni’. \textit{Pamiętnik Literacki} 96/2 (2005): 13–15.
\lit Pilat, R. \textit{Pieśń ‘Bogarodzica’. I. Restytucyja}. Kraków, 1879.
\lit Surzyński, J. \textit{Polskie pieśni Kościoła katolickiego od
  najdawniejszych czasów do końca XVI stulecia}. Poznań, 1891.
\lit Woronczak, J. ‘Tropy i sekwencje w literaturze polskiej do połowy XVI
  wieku’. \textit{Pamiętnik Literacki} 43 (1952): 353–54.
\lit Woronczak, J., ed., with E. Ostrowska and H. Feicht.
  \textit{Bogurodzica}. Biblioteka Pisarzy Polskich A~1. Wrocław:
  Ossolineum, 1962.
\end{literature}

\end{document}
